Regularization is central to training neural networks that generalize well. This report empirically compares common regularization approaches---including $L_2$ weight decay, $L_1$ penalty, dropout, and early stopping---with a focus on how they affect convergence speed and generalization, and contrasts them against \emph{synaptic neural balance} (neural balancing), a weight re-scaling transformation that preserves the network function while reducing a chosen weight cost.

We evaluate all methods under a controlled protocol on MNIST with a small fully connected network ($784 \to 256 \to 10$, ReLU), using GP Bayesian optimization to select hyperparameters on a held-out validation split and reporting test metrics across \NumSeeds{} fixed random seeds. Classical regularizers ($L_1$, $L_2$, early stopping) all converge to approximately $91.96\%$ test accuracy---slightly below the plain baseline ($92.29\%$)---while dropout reaches $\RegMnistDropoutFinalTestAccPct\%$. By contrast, synaptic balance with an $L_1$ cost achieves $\RegMnistSynBalLOneFinalTestAccPct\%$ test accuracy, $\RegMnistSynBalLOneGainOverBestNonBalancePpts$ percentage points above the best non-balancing baseline, and converges to $\TargetTauPct\%$ accuracy in just $\RegMnistSynBalLOneEpochsAtTau$ epochs compared to $67$--$70$ epochs for classical methods. Synaptic balance with $L_2$ cost also improves on the best non-balancing baseline by $\RegMnistSynBalLTwoGainOverBestNonBalancePpts$ percentage points with faster convergence ($\RegMnistSynBalLTwoEpochsAtTau$ epochs). These results position neural balancing---particularly with $L_1$ cost---as a complementary tool that provides measurable benefits in both optimization speed and generalization in this controlled setting.